\documentclass[a4paper,12pt]{article}

\usepackage[T2A]{fontenc}
\usepackage[utf8]{inputenc}
\usepackage[russian]{babel}
\usepackage{amsmath,amssymb,mathtools}
\usepackage{PTSans}
\renewcommand{\familydefault}{\sfdefault}
\usepackage{icomma}
\usepackage[a4paper,margin=2.1cm,headheight=15pt]{geometry}
\usepackage{microtype}
\microtypesetup{nopatch=footnote}
\usepackage{tikz}
\usetikzlibrary{arrows.meta,positioning}
\usepackage{tabularx,booktabs}
\usepackage{enumitem}
\usepackage{hyperref}
\usepackage{parskip}
\usepackage{fancyhdr}
\usepackage{setspace}
\usepackage{xcolor}

\definecolor{accent}{HTML}{008080}
\definecolor{darkaccent}{HTML}{004D4D}
\definecolor{textgray}{HTML}{333333}
\definecolor{softgray}{HTML}{F3F5F5}
\definecolor{linegray}{HTML}{A9B4B4}

\hypersetup{
  unicode=true,
  colorlinks=true,
  linkcolor=darkaccent,
  urlcolor=darkaccent,
  pdfauthor={Лёвин Артём Александрович},
  pdftitle={Чек-лист: задачи на производительность и совместную работу}
}

\geometry{includeheadfoot}
\onehalfspacing
\setlength{\parindent}{0pt}
\setlist[itemize]{leftmargin=1.5em,itemsep=0.2em,topsep=0.25em}
\setlist[enumerate]{leftmargin=1.7em,itemsep=0.45em,topsep=0.3em}
\renewcommand{\arraystretch}{1.3}

\pagestyle{fancy}
\fancyhf{}
\fancyhead[L]{\small\textcolor{textgray}{Задачи на производительность}}
\fancyhead[R]{\small\textcolor{textgray}{Кирилл Зиновьев · 7 класс}}
\fancyfoot[C]{\small\textcolor{textgray}{\thepage}}
\renewcommand{\headrulewidth}{0.35pt}
\renewcommand{\footrulewidth}{0pt}

\newcommand{\lessonsection}[1]{%
  \vspace{0.35em}
  {\Large\bfseries\textcolor{darkaccent}{#1}}\par
  \vspace{0.2em}
  {\color{accent}\rule{\linewidth}{0.6pt}}\par
  \vspace{0.4em}
}

\newcommand{\keyidea}[1]{\textbf{\textcolor{darkaccent}{#1}}}
\newcommand{\unit}[1]{\,\text{#1}}

\begin{document}

\hypersetup{pageanchor=false}
\begin{titlepage}
\thispagestyle{empty}
\centering
\vspace*{3.0cm}

{\Large\textcolor{accent}{\textbf{ЧЕК-ЛИСТ}}\par}
\vspace{1.0cm}

{\Huge\bfseries\textcolor{darkaccent}{Задачи на производительность}}\par
\vspace{0.25cm}
{\Huge\bfseries\textcolor{darkaccent}{и совместную работу}}\par

\vspace{1.4cm}

{\large 7 класс\par}

\vspace{2.0cm}

{\large\textbf{Лёвин Артём Александрович}\par}
\vspace{0.35cm}
{\large эксклюзивно для Кирилла Зиновьева\par}

\vfill

{\large \today\par}

\vspace{2.0cm}

\begin{tikzpicture}[remember picture,overlay]
  \draw[accent!45,line width=1.2pt]
    ([xshift=1.2cm,yshift=1.25cm]current page.south west)
    .. controls
    ([xshift=5.1cm,yshift=2.05cm]current page.south west) and
    ([xshift=-6.0cm,yshift=0.55cm]current page.south east)
    ..
    ([xshift=-1.2cm,yshift=1.45cm]current page.south east);

  \draw[darkaccent!20,line width=0.6pt]
    ([xshift=1.2cm,yshift=0.95cm]current page.south west)
    .. controls
    ([xshift=6.0cm,yshift=1.65cm]current page.south west) and
    ([xshift=-5.2cm,yshift=0.35cm]current page.south east)
    ..
    ([xshift=-1.2cm,yshift=1.10cm]current page.south east);
\end{tikzpicture}

\end{titlepage}

\hypersetup{pageanchor=true}
\setcounter{page}{1}

\lessonsection{1. Главная модель}

В задачах о работе удобно держать в голове три величины:

\begin{center}
\begin{tabularx}{0.92\linewidth}{@{}lXl@{}}
\toprule
\textbf{Величина} & \textbf{Смысл} & \textbf{Единицы в задачах урока} \\
\midrule
Производительность $p$ &
сколько деталей выполняется за единицу времени &
деталей/ч \\
Время $t$ &
сколько времени выполнялась работа &
ч \\
Результат $A$ &
сколько деталей изготовлено &
деталей \\
\bottomrule
\end{tabularx}
\end{center}

\keyidea{Основная формула:}
\[
\boxed{A=pt}
\qquad\Longleftrightarrow\qquad
\boxed{t=\frac{A}{p}}
\qquad\Longleftrightarrow\qquad
\boxed{p=\frac{A}{t}}.
\]

Мнемоника: \textbf{П $\cdot$ В $=$ Р} ---
производительность, умноженная на время, даёт результат.

\begin{center}
\begin{tikzpicture}[font=\small]
  \node[
    draw=darkaccent,
    rounded corners=2pt,
    line width=0.8pt,
    minimum width=2.3cm,
    minimum height=0.9cm
  ] (p) {$p$ --- производительность};

  \node[
    draw=darkaccent,
    rounded corners=2pt,
    line width=0.8pt,
    minimum width=1.8cm,
    minimum height=0.9cm,
    right=1.2cm of p
  ] (t) {$t$ --- время};

  \node[
    draw=accent,
    rounded corners=2pt,
    line width=1pt,
    minimum width=2.0cm,
    minimum height=0.9cm,
    right=1.2cm of t
  ] (a) {$A$ --- результат};

  \draw[-{Latex[length=2.5mm]},line width=0.9pt,textgray]
    (p) -- node[above] {$\times$} (t);

  \draw[-{Latex[length=2.5mm]},line width=0.9pt,textgray]
    (t) -- (a);
\end{tikzpicture}
\end{center}

\textbf{Ограничение.}
В формулах с делением знаменатель не должен быть равен нулю.
В задаче о производительности $p>0$, а при вычислении
\[
p=\frac{A}{t}
\]
также требуется $t>0$.

\lessonsection{2. Совместная работа}

Если два рабочих выполняют одну работу одновременно,
их производительности \textbf{складываются}:
\[
p_{\Sigma}=p_1+p_2.
\]

Следовательно, за одинаковое время $t$ они выполнят
\[
A=(p_1+p_2)t.
\]

\begin{center}
\begin{tabularx}{0.92\linewidth}{@{}Xccc@{}}
\toprule
& \textbf{1-й рабочий} & \textbf{2-й рабочий} & \textbf{Вместе} \\
\midrule
Производительность & $p_1$ & $p_2$ & $p_1+p_2$ \\
Время & $t$ & $t$ & $t$ \\
Результат & $p_1t$ & $p_2t$ & $(p_1+p_2)t$ \\
\bottomrule
\end{tabularx}
\end{center}

\keyidea{Важно:}
складываются именно производительности.
Время не складывают, если рабочие трудятся \textbf{одновременно}.

\lessonsection{3. Четыре базовых типа задач}

\textbf{Тип 1. Искомое --- время.}

Один рабочий делает $7$ деталей/ч, второй --- $9$ деталей/ч.
Вместе нужно изготовить $64$ детали.

\[
(7+9)t=64,
\qquad
16t=64,
\qquad
\boxed{t=4\unit{ч}}.
\]

\textbf{Тип 2. Искомое --- производительность второго.}

Первый делает $7$ деталей/ч.
Вместе за $4$ ч изготовлено $64$ детали.

\[
(7+p_2)\cdot4=64,
\qquad
7+p_2=16,
\qquad
\boxed{p_2=9\unit{деталей/ч}}.
\]

\textbf{Тип 3. Производительности связаны разностью.}

Второй делает на $2$ детали/ч больше первого.
Вместе за $4$ ч изготовлено $64$ детали.

Пусть
\[
p_1=x,
\qquad
p_2=x+2.
\]

Тогда
\[
(x+x+2)\cdot4=64,
\]
\[
2x+2=16,
\qquad
2x=14,
\qquad
x=7.
\]

Значит,
\[
\boxed{p_1=7\unit{деталей/ч}},
\qquad
\boxed{p_2=9\unit{деталей/ч}}.
\]

\textbf{Тип 4. Разность результатов за одно и то же время.}

Первый делает $7$ деталей/ч, второй --- $9$ деталей/ч.
Через сколько часов второй изготовит на $8$ деталей больше первого?

За $t$ часов первый изготовит $7t$ деталей,
а второй --- $9t$ деталей. Поэтому
\[
9t-7t=8,
\]
\[
2t=8,
\qquad
\boxed{t=4\unit{ч}}.
\]

Тот же смысл можно записать иначе.
Если первый изготовил $x$ деталей, то второй изготовил $x+8$ деталей.
Так как они работали одинаковое время,
\[
\frac{x}{7}=\frac{x+8}{9}.
\]

После решения получаем $x=28$, а значит
\[
t=\frac{28}{7}=4\unit{ч}.
\]

\lessonsection{4. Как переводить текст задачи в уравнение}

\begin{enumerate}
  \item
  \textbf{Определите, что известно:}
  производительность, время, результат или связь между величинами.

  \item
  \textbf{Выберите неизвестное.}
  Удобно обозначать через $x$ либо величину, о которой спрашивают,
  либо производительность первого рабочего.

  \item
  \textbf{Заполните таблицу} $p$--$t$--$A$.
  Недостающие клетки восстанавливайте по формуле
  \[
  A=pt.
  \]

  \item
  \textbf{Найдите фразу-связку:}
  «вместе», «на $k$ больше», «за одинаковое время»,
  «начал раньше», «работал отдельно».

  \item
  \textbf{Составьте уравнение}
  по общему результату, равенству результатов
  или их разности.

  \item
  \textbf{Решите уравнение и вернитесь к условию.}
  Укажите единицы измерения и запишите полный ответ.
\end{enumerate}

\begin{center}
\begin{tikzpicture}[node distance=7mm and 9mm,font=\small]
\tikzset{
  algstep/.style={
    draw=darkaccent,
    rounded corners=2pt,
    minimum width=3.0cm,
    minimum height=0.8cm,
    align=center
  }
}

\node[algstep] (s1) {Известные величины};
\node[algstep,right=of s1] (s2) {Таблица $p$--$t$--$A$};
\node[algstep,right=of s2] (s3) {Уравнение};
\node[algstep,below=of s3] (s4) {Решение};
\node[algstep,left=of s4] (s5) {Проверка единиц};
\node[algstep,left=of s5] (s6) {Ответ};

\draw[-{Latex[length=2.5mm]},line width=0.8pt,accent]
  (s1) -- (s2);
\draw[-{Latex[length=2.5mm]},line width=0.8pt,accent]
  (s2) -- (s3);
\draw[-{Latex[length=2.5mm]},line width=0.8pt,accent]
  (s3) -- (s4);
\draw[-{Latex[length=2.5mm]},line width=0.8pt,accent]
  (s4) -- (s5);
\draw[-{Latex[length=2.5mm]},line width=0.8pt,accent]
  (s5) -- (s6);
\end{tikzpicture}
\end{center}

\lessonsection{5. Слова, которые меняют модель}

\begin{center}
\begin{tabularx}{\linewidth}{@{}p{0.30\linewidth}X@{}}
\toprule
\textbf{Фраза в условии} & \textbf{Математическая запись} \\
\midrule

«Второй делает на $d$ деталей/ч больше»
&
$p_2=p_1+d$
\\

«Второй делает на $d$ деталей/ч меньше»
&
$p_2=p_1-d$
\\

«Второй изготовил на $D$ деталей больше»
при одинаковом времени
&
$A_2=A_1+D$,
либо
$(p_2-p_1)t=D$
\\

«Работали вместе $t$ часов»
&
$A=(p_1+p_2)t$
\\

«Первый начал на $h$ часов раньше»
&
у первого время $t+h$,
у второго --- $t$
\\

«Первый работал отдельно, затем вместе»
&
результаты отдельных этапов складываются
\\

\bottomrule
\end{tabularx}
\end{center}

\lessonsection{6. Типичные ошибки и быстрая самопроверка}

\begin{itemize}

  \item
  \textbf{Смешение величин.}
  Нельзя складывать
  $7\unit{деталей/ч}$ и $4\unit{ч}$.
  В уравнении должны сочетаться совместимые величины.

  \item
  \textbf{Неверная совместная производительность.}
  Для двух рабочих
  \[
  p_{\Sigma}=p_1+p_2,
  \]
  а не $p_1p_2$.

  \item
  \textbf{Путаница «на» и «в».}
  «На $3$ больше» означает
  \[
  x+3,
  \]
  а «в $3$ раза больше» означает
  \[
  3x.
  \]

  \item
  \textbf{Одинаковое время используется не всегда.}
  Если один рабочий начал раньше или закончил позже,
  времена нужно обозначить отдельно.

  \item
  \textbf{Потерян множитель времени.}
  Если известен общий результат,
  уравнение обычно содержит произведение
  \[
  (p_1+p_2)t.
  \]

  \item
  \textbf{Не проверен смысл ответа.}
  Производительность должна быть положительной,
  время --- неотрицательным,
  а подстановка найденных значений должна
  восстановить исходный результат.

\end{itemize}

\textbf{Контрольный вопрос перед ответом:}
полученная величина имеет правильные единицы ---
часы, детали или детали/ч?

\clearpage

\lessonsection{Тренировочный блок}

\textbf{Средний уровень}

\begin{enumerate}[label=\textbf{\arabic*.}]

  \item
  Один рабочий изготовляет $6$ деталей в час,
  другой --- $8$.
  За сколько часов они изготовят вместе $70$ деталей?

  \item
  Первый рабочий изготовляет $7$ деталей в час.
  Работая вместе со вторым,
  они за $5$ ч изготовили $90$ деталей.
  Сколько деталей в час изготовляет второй рабочий?

  \item
  Второй рабочий изготовляет за час
  на $4$ детали больше, чем первый.
  Вместе за $6$ ч они изготовили $120$ деталей.
  Найдите производительность каждого.

\end{enumerate}

\vspace{0.5em}

\textbf{Выше среднего}

\begin{enumerate}[label=\textbf{\arabic*.},start=4]

  \item
  Первый рабочий изготовляет $7$ деталей в час,
  второй --- $10$.
  Через сколько часов второй изготовит
  на $18$ деталей больше первого,
  если они начнут одновременно?

  \item
  Первый рабочий изготовляет $12$ деталей в час,
  второй --- $15$.
  Через сколько часов второй изготовит
  на $21$ деталь больше первого,
  если время их работы одинаково?

  \item
  Первый рабочий изготовляет $6$ деталей в час.
  Сначала оба рабочих трудились вместе $3$ ч,
  затем первый работал один ещё $2$ ч.
  Всего изготовили $72$ детали.
  Найдите производительность второго рабочего.

  \item
  Первый рабочий изготовляет $8$ деталей в час.
  Он работал один $2$ ч,
  затем к нему присоединился второй,
  и вместе они работали ещё $4$ ч.
  Всего изготовили $100$ деталей.
  Сколько деталей в час изготовляет второй рабочий?

  \item
  Первый рабочий изготовляет $9$ деталей в час,
  второй --- $11$.
  Первый начал работу на $2$ ч раньше.
  После начала работы второго общий результат,
  включая уже выполненную первым работу,
  составил $138$ деталей.
  Сколько часов рабочие трудились вместе?

  \item
  Второй рабочий изготовляет
  на $3$ детали в час больше первого.
  Первый работал $3$ ч,
  второй --- $2$ ч.
  Всего они изготовили $46$ деталей.
  Найдите производительность каждого.

\end{enumerate}

\vspace{0.5em}

\textbf{Сложная задача}

\begin{enumerate}[label=\textbf{\arabic*.},start=10]

  \item
  Второй рабочий изготовляет
  на $5$ деталей в час больше первого.
  За $4$ ч работы первого
  и $3$ ч работы второго
  они изготовили всего $71$ деталь.

  Найдите производительность каждого рабочего.

  Затем определите,
  на сколько часов раньше второй рабочий
  изготовит $104$ детали,
  если каждый будет работать отдельно
  с найденной производительностью.

\end{enumerate}

\clearpage

\lessonsection{Ответы к тренировочному блоку}

\begin{table}[ht]
\centering
\caption{Краткие ответы}

\begin{tabularx}{0.92\linewidth}{@{}cX@{}}
\toprule
\textbf{№} & \textbf{Ответ} \\
\midrule

1 & $5$ ч. \\

2 & $11$ деталей/ч. \\

3 & $8$ деталей/ч и $12$ деталей/ч. \\

4 & $6$ ч. \\

5 & $7$ ч. \\

6 & $14$ деталей/ч. \\

7 & $13$ деталей/ч. \\

8 & $6$ ч совместной работы. \\

9 & $8$ деталей/ч и $11$ деталей/ч. \\

10 &
$8$ деталей/ч и $13$ деталей/ч;
второй изготовит $104$ детали на $5$ ч раньше.
\\

\bottomrule
\end{tabularx}
\end{table}

\vspace{1em}

\textbf{Финальная проверка.}
Если ответ найден верно,
подстановка в исходное условие восстанавливает
все данные: производительности, время
и общий результат.

\end{document}